% JopeDime: where our evolved models win. A plain-language brief.
% Every number and chart comes from winners-data/ (build-winners.sh), computed from cross-validated
% evidence on training data only (5 folds of train.csv; test.csv is never read). The full technical
% record of all 14 cases is campaign-technical.tex. This is not the white paper (docs/whitepaper/).
\documentclass[11pt]{article}
\usepackage[a4paper,margin=20mm,headheight=14pt]{geometry}
\usepackage[T1]{fontenc}
\usepackage[utf8]{inputenc}
\usepackage{lmodern}
\usepackage[scaled=0.92]{helvet}
\renewcommand{\familydefault}{\sfdefault}
\usepackage{microtype}
\usepackage[table]{xcolor}
\usepackage{booktabs,array,tabularx}
\usepackage{siunitx}
\sisetup{mode=text,round-mode=figures,round-precision=3,exponent-product=\times,group-digits=integer,group-separator={,},group-minimum-digits=4}
\usepackage{enumitem}
\setlist{nosep,leftmargin=1.3em,itemsep=3pt}
\usepackage{pgfplots,pgfplotstable}
\pgfplotsset{compat=1.18}
\usetikzlibrary{arrows.meta,positioning}
\usepackage[most]{tcolorbox}
\usepackage{fancyhdr}
\usepackage{titlesec}
\usepackage[hidelinks]{hyperref}

\definecolor{ours}{HTML}{1F5FAD}
\definecolor{oursmid}{HTML}{5B8BD0}
\definecolor{oursstart}{HTML}{B9CDE8}
\definecolor{eruns}{HTML}{7FB6A6}
\definecolor{alt}{HTML}{8C8C8C}
\definecolor{bad}{HTML}{C0504D}
\definecolor{win}{HTML}{0B8A6F}
\definecolor{ink}{HTML}{1D1D1F}
\definecolor{panel}{HTML}{F3F5F8}

\titleformat{\section}{\Large\bfseries\color{ink}}{}{0em}{}
\titleformat{\subsection}{\large\bfseries\color{ours}}{}{0em}{}
\titlespacing*{\section}{0pt}{12pt}{5pt}
\titlespacing*{\subsection}{0pt}{9pt}{3pt}
\pagestyle{fancy}\fancyhf{}
\fancyhead[L]{\small\color{alt}JopeDime --- where our evolved models win}
\fancyhead[R]{\small\color{alt}data of \BriefDataStamp}
\fancyfoot[C]{\small\color{alt}\thepage}
\renewcommand{\headrulewidth}{0.3pt}
\setlength{\parskip}{4pt}\setlength{\parindent}{0pt}

\newtcolorbox{plainbox}[1]{colback=panel,colframe=ours,boxrule=0.6pt,arc=3pt,left=8pt,right=8pt,top=5pt,bottom=5pt,fonttitle=\bfseries,title=#1}
\newtcolorbox{resultbox}{colback=win!6,colframe=win,boxrule=0.9pt,arc=3pt,left=9pt,right=9pt,top=7pt,bottom=7pt}
\newtcolorbox{carebox}{colback=white,colframe=alt,boxrule=0.5pt,arc=3pt,left=8pt,right=8pt,top=5pt,bottom=5pt,fonttitle=\bfseries\color{ink},coltitle=ink,colbacktitle=panel,title=Keep in mind}

\newcommand{\WinTwoChainCols}{A7,A11,A13,E7,E11,E13}
\newcommand{\WinTwoBetter}{24}
\newcommand{\WinTwoBestBetter}{27}
\newcommand{\WinTwoOursMedianRmse}{0.01324}
\newcommand{\WinTwoArmEnds}{E: 0.445, 0.438, 0.437; A: 0.518, 0.494, 0.472}
\newcommand{\WinTwoArmStarts}{E: 0.492, 0.493, 0.491; A: 0.778, 0.758, 0.687}
\newcommand{\WinTwoName}{Hourly Day-Ahead Electricity Prices with Exogenous Variables}
\newcommand{\WinTwoRows}{6720}
\newcommand{\WinTwoFeatures}{11}
\newcommand{\WinTwoOurs}{0.458532}
\newcommand{\WinTwoOursBest}{0.43698}
\newcommand{\WinTwoChains}{6}
\newcommand{\WinTwoStart}{0.589897}
\newcommand{\WinTwoAlt}{0.600557}
\newcommand{\WinTwoAltModel}{xgboost\_default}
\newcommand{\WinTwoAltCount}{14}
\newcommand{\WinTwoMargin}{-23.6}
\newcommand{\WinTwoAltAcc}{0.013226}
\newcommand{\WinTwoAltCpu}{0.383}
\newcommand{\WinTwoAltFull}{14}
\newcommand{\WinTwoAltPartial}{3}
\newcommand{\WinTwoGens}{112}
\newcommand{\WinTwoEvals}{20802}
\newcommand{\WinTwoStartPlain}{boosted decision trees}
\newcommand{\WinTwoStartProgram}{(pipe X (histgbm bins=255 depth=6 lambda=1 lr=0.3 min\_leaf=1 rounds=100))}
\newcommand{\WinTwoStartObj}{0.491238}
\newcommand{\WinTwoStartRmse}{0.014172}
\newcommand{\WinTwoStartFitCpuSeconds}{0.07499291380003412}
\newcommand{\WinTwoStartFitSeconds}{0.0752334622}
\newcommand{\WinTwoStartRowLatencyPninetynineUs}{n/a}
\newcommand{\WinTwoStartModelBytes}{441716.2}
\newcommand{\WinTwoStartProgramNodes}{1}
\newcommand{\WinTwoStartPredictUsPerRow}{1.987201122522362}
\newcommand{\WinTwoStartGen}{0}
\newcommand{\WinTwoStartModelKB}{442}
\newcommand{\WinTwoEndPlain}{boosted decision trees}
\newcommand{\WinTwoEndProgram}{(pipe X (histgbm bins=256 depth=4 lambda=1 lr=0.3 min\_leaf=1 rounds=100))}
\newcommand{\WinTwoEndObj}{0.43698}
\newcommand{\WinTwoEndRmse}{0.012873}
\newcommand{\WinTwoEndFitCpuSeconds}{0.04779845899997515}
\newcommand{\WinTwoEndFitSeconds}{0.047974075600000006}
\newcommand{\WinTwoEndRowLatencyPninetynineUs}{n/a}
\newcommand{\WinTwoEndModelBytes}{139062.2}
\newcommand{\WinTwoEndProgramNodes}{1}
\newcommand{\WinTwoEndPredictUsPerRow}{1.0751693641099809}
\newcommand{\WinTwoEndGen}{26}
\newcommand{\WinTwoEndModelKB}{139}
\newcommand{\WinTwoBestChain}{E13}
\newcommand{\WinTwoChanges}{the same kind of model (boosted decision trees) with new settings: histogram bins 255 $\to$ 256, tree depth 6 $\to$ 4}
\newcommand{\WinTwoCpuRatio}{8}
\newcommand{\WinTwoFoldWins}{5}
\newcommand{\WinTwoPartialNote}{AutoGluon (fold 1): 379.3 against our 0.2869 on the same fold; Kumo-Tabular (fold 1): 0.8951 against our 0.2869 on the same fold; TabPFN-v2 (fold 1): 0.3921 against our 0.2869 on the same fold}
\newcommand{\WinTwoPartialList}{AutoGluon 379, Kumo-Tabular 0.895, TabPFN-v2 0.392}
\newcommand{\WinTwoOursFoldOneRmse}{0.008191}
\newcommand{\WinTwoOursFoldOneObj}{0.2869}
\newcommand{\WinTwoAutoGluonFoldOneAcc}{0.007457}
\newcommand{\WinTwoAutoGluonFoldOneCpu}{833.9}
\newcommand{\WinTwoAutoGluonFoldOneObj}{379.3}
\newcommand{\WinTwoAutoGluonMean}{}
\newcommand{\WinTwoTabPFNFoldOneAcc}{0.009076}
\newcommand{\WinTwoTabPFNFoldOneCpu}{0.2186}
\newcommand{\WinTwoTabPFNFoldOneObj}{0.3921}
\newcommand{\WinTwoTabPFNMean}{}
\newcommand{\WinTwoKumoFoldOneAcc}{0.02457}
\newcommand{\WinTwoKumoFoldOneCpu}{0.2253}
\newcommand{\WinTwoKumoFoldOneObj}{0.8951}
\newcommand{\WinTwoKumoMean}{}
\newcommand{\WinTwoXGBTunedFoldOneAcc}{0.0119}
\newcommand{\WinTwoXGBTunedFoldOneCpu}{0.7002}
\newcommand{\WinTwoXGBTunedFoldOneObj}{0.7022}
\newcommand{\WinTwoXGBTunedMean}{1.296}
\newcommand{\WinTwoLGBTunedFoldOneAcc}{0.01151}
\newcommand{\WinTwoLGBTunedFoldOneCpu}{1.364}
\newcommand{\WinTwoLGBTunedFoldOneObj}{0.9911}
\newcommand{\WinTwoLGBTunedMean}{1.269}
\newcommand{\WinTwoBenchMax}{1.621}
\newcommand{\WinTwoBenchRows}{17}
\newcommand{\WinTwoBenchLast}{16}
\newcommand{\BriefCvStamp}{20261103T170155Z-cv5}
\newcommand{\BriefDataStamp}{2026-11-03T19:20:45Z}
\newcommand{\BriefInputCount}{225}
\newcommand{\BriefInputsSha}{e5504c3944eafd88}

\newcommand{\W}[2]{\csname Win#1#2\endcsname}   % \W{Two}{Ours}
\newcommand{\bignum}[2]{\begin{minipage}[t]{0.31\linewidth}\centering{\fontsize{22}{24}\selectfont\bfseries #1}\\[2pt]{\small #2}\end{minipage}}

% Benchmark chart: every alternative's cross-validated score, our start, typical and best runs.
\newcommand{\benchchart}[3]{% #1 case number, #2 case word, #3 x label
  \pgfplotstableread[col sep=comma]{winners-data/#1-bench.csv}\benchtable
  \begin{tikzpicture}\begin{axis}[xbar,width=0.74\linewidth,y=0.44cm,bar width=0.30cm,bar shift=0pt,
      enlarge y limits={abs=0.45},xmin=0,xmax=\W{#2}{BenchMax},ytick={0,...,\W{#2}{BenchLast}},yticklabels from table={\benchtable}{label},
      y tick label style={font=\small},x tick label style={font=\footnotesize,/pgf/number format/fixed},
      xlabel={#3},label style={font=\small},axis x line*=bottom,axis y line*=left,unbounded coords=discard,
      nodes near coords,point meta=explicit symbolic,every node near coord/.style={font=\scriptsize,anchor=west,color=ink}]
    \addplot[fill=alt!50,draw=none] table[x=alt,y expr=\coordindex,meta=text]{\benchtable};
    \addplot[fill=bad!30,draw=bad] table[x=bad,y expr=\coordindex,meta=text]{\benchtable};
    \addplot[fill=oursstart,draw=none] table[x=start,y expr=\coordindex,meta=text]{\benchtable};
    \addplot[fill=oursmid,draw=none] table[x=ours,y expr=\coordindex,meta=text]{\benchtable};
    \addplot[fill=ours,draw=none] table[x=best,y expr=\coordindex,meta=text]{\benchtable};
  \end{axis}\end{tikzpicture}}

% Evolution chart: the best cross-validated score so far, per run; the best run is bold.
\newcommand{\evochart}[6]{% #1 case number, #2 case word, #3 csv suffix, #4 y label, #5 reference value, #6 reference label
  \begin{tikzpicture}\begin{axis}[width=0.92\linewidth,height=6.6cm,grid=major,grid style={gray!18},unbounded coords=jump,
      xlabel={generations of evolution (across all rounds)},ylabel={#4},label style={font=\small},tick label style={font=\footnotesize},
      scaled y ticks=false,y tick label style={/pgf/number format/fixed,/pgf/number format/precision=2},xmin=0,
      legend style={font=\footnotesize,draw=none,fill=white,fill opacity=0.9,text opacity=1,at={(0.99,0.99)},anchor=north east,cells={anchor=west}}]
    \foreach \c in {E7,E11,E13} {\addplot[eruns,thick,const plot,forget plot] table[col sep=comma,x=g,y=\c]{winners-data/#1-#3.csv};}
    \foreach \c in {A7,A11,A13} {\addplot[oursstart,thick,const plot,forget plot] table[col sep=comma,x=g,y=\c]{winners-data/#1-#3.csv};}
    \addplot[ours,ultra thick,const plot] table[col sep=comma,x=g,y=\W{#2}{BestChain}]{winners-data/#1-#3.csv};\addlegendentry{best run}
    \addlegendimage{eruns,thick}\addlegendentry{runs started from industry-standard settings}
    \addlegendimage{oursstart,thick}\addlegendentry{runs started from our neural model's suggestions}
    \addplot[bad,dashed,thick,domain=0:\W{#2}{Gens}] {#5};\addlegendentry{#6}
    \addplot[only marks,mark=*,mark size=3pt,color=ours] coordinates {(0,\W{#2}{StartObj}) (\W{#2}{EndGen},\W{#2}{EndObj})};
    \node[font=\footnotesize\bfseries,color=ours,anchor=south west] at (axis cs:1,\W{#2}{StartObj}) {start \num{\W{#2}{StartObj}}};
    \node[font=\footnotesize\bfseries,color=ours,anchor=north west] at (axis cs:\W{#2}{EndGen},\W{#2}{EndObj}) {end \num{\W{#2}{EndObj}}};
  \end{axis}\end{tikzpicture}}

\begin{document}

\begin{center}
{\LARGE\bfseries\color{ink} Where JopeDime's evolved models win}\\[6pt]
{\large\color{ours} Prediction programs tuned to what the business actually pays for}\\[5pt]
{\small\color{alt} Case-study brief \textbullet{} every score cross-validated on training data \textbullet{} \BriefDataStamp}
\end{center}

\section{The short version}
Most machine-learning contests reward accuracy and nothing else. A business pays for more than accuracy:
for the computer time it takes to retrain a model, for how fast it answers, and for how much it costs to run.
JopeDime lets an evolutionary search write prediction programs and judges every candidate on a
\textbf{business scorecard} that includes those costs.

We ran 14 business cases. In each, we compared our evolved programs with well-known alternatives: in this case, \W{Two}{AltFull} of them on all
5 folds (including XGBoost, LightGBM, CatBoost, tuned versions of the most popular libraries, and classic statistical methods) and
\W{Two}{AltPartial} slower ones (AutoGluon, TabPFN and Kumo) on one fold.
We count a case only if we win it \emph{clearly}: no alternative comes close on that case's scorecard.
\textbf{One case passes that test}, and this brief is about it.

\begin{resultbox}
{\large\bfseries\color{win} Hourly electricity-price forecasting, with every retrain billed by CPU time}\par\medskip
\bignum{\W{Two}{Better}\,\%}{better business score than the best alternative (XGBoost), for our typical run}\hfill
\bignum{\W{Two}{BestBetter}\,\%}{better for our best run, which wins on \W{Two}{FoldWins} of 5 data folds}\hfill
\bignum{\W{Two}{CpuRatio}\,\texttimes}{less training CPU per refit than XGBoost (\num[round-precision=2]{\W{Two}{AltCpu}}\,s vs \num[round-precision=2]{\W{Two}{EndFitCpuSeconds}}\,s), for our best program}
\end{resultbox}

\begin{plainbox}{How we kept it fair}
\begin{itemize}
\item \textbf{The same data and the same scorecard for everyone.} Every model, ours and the alternatives, is fitted on the same rows and scored
by the same program.
\item \textbf{Cross-validation on the training data only.} The training data is split into 5 parts (``folds''). Each model is
trained on 4 parts and scored on the fifth, five times over, and the score is the average. The separate test data was never used.
\item \textbf{Every number comes from the run records.} The brief is regenerated from them by a script; nothing is typed by hand.
\end{itemize}
\end{plainbox}

\newpage
\section{The case: forecasting tomorrow's hourly electricity prices}
\subsection{The question}
An energy trader re-estimates a price-forecast model every hour, across many markets, and every refit is paid
for in computer time. \emph{Which model gives the best forecasts once you also pay for the computer time it takes to retrain?}
The data are \num{\W{Two}{Rows}} hours of day-ahead electricity prices with \W{Two}{Features} inputs (the market, the time, the day of the week
and two outside variables), from a public benchmark collection (Nixtla).

\subsection{The scorecard}
\textbf{Score = forecast error + a charge for training time.} The error is measured relative to how much prices normally move.
The charge adds a tenth of a typical price swing for every 0.22 seconds of CPU time a retrain uses.
A model that forecasts equally well but retrains in a fifth of the time therefore gets the better score.
\textbf{Lower is better.}

\subsection{How everyone scored}
\begin{center}
\benchchart{02}{Two}{business score, averaged over 5 folds (lower is better)}
\end{center}
{\small Grey: \W{Two}{AltFull} alternatives, each scored on all 5 folds (``tuned'' versions searched 40 settings inside every training fold).
Light blue: where our best run \emph{started}, from our engine's version of XGBoost's default settings.
Blue: our typical run (the middle of 6 runs). Dark blue: our best run. Very large scores are cut at the edge of the chart, with their value written beside them.}

\smallskip
{\small \W{Two}{AltPartial} slower alternatives were scored on one fold only, so they are not in the chart. On that fold they scored
\W{Two}{PartialList}, against \num{\W{Two}{OursFoldOneObj}} for our best program on the same fold.}

\newpage
\subsection{How evolution improved the score}
\begin{center}
\evochart{02}{Two}{evo}{business score of the best program so far}{\W{Two}{Alt}}{best alternative: XGBoost (\num{\W{Two}{Alt}})}
\end{center}
{\small Each line is one run of the search: six runs in all (two ways of choosing the starting programs $\times$ three random seeds), each of up to six rounds.
A line drops each time the run finds a better program, scored over all 5 folds. Together the runs evaluated about
\num{\W{Two}{Evals}} programs.}

\subsection{Where the best run started, and where it ended up}
\noindent\begin{tcolorbox}[enhanced,sidebyside,sidebyside align=top seam,colback=white,colframe=ours!60,boxrule=0.6pt,arc=3pt,
    lefthand width=0.47\linewidth,left=7pt,right=7pt,top=6pt,bottom=6pt,segmentation style={dashed,ours!40}]
  {\bfseries\color{alt}Where it started} {\small(generation 0)}\par\smallskip
  Boosted decision trees with XGBoost's standard settings, as our engine runs them.\par\smallskip
  \begin{tabular}{@{}ll@{}}
  business score & \textbf{\num{\W{Two}{StartObj}}}\\
  forecast error (RMSE) & \num{\W{Two}{StartRmse}}\\
  training CPU per refit & \num[round-precision=2]{\W{Two}{StartFitCpuSeconds}}\,s\\
  model size & \W{Two}{StartModelKB}\,KB
  \end{tabular}\par\smallskip
  {\scriptsize\ttfamily\color{alt}\W{Two}{StartProgram}}
\tcblower
  {\bfseries\color{ours}Where it ended up} {\small(found in generation \W{Two}{EndGen})}\par\smallskip
  What changed: \W{Two}{Changes}.\par\smallskip
  \begin{tabular}{@{}ll@{}}
  business score & \textbf{\num{\W{Two}{EndObj}}}\\
  forecast error (RMSE) & \num{\W{Two}{EndRmse}}\\
  training CPU per refit & \num[round-precision=2]{\W{Two}{EndFitCpuSeconds}}\,s\\
  model size & \W{Two}{EndModelKB}\,KB
  \end{tabular}\par\smallskip
  {\scriptsize\ttfamily\color{alt}\W{Two}{EndProgram}}
\end{tcolorbox}

\subsection{What this means}
\begin{itemize}
\item \textbf{More accurate and cheaper at the same time.} Our best program's forecast error (\num{\W{Two}{EndRmse}}) is lower than that of every
alternative scored on all 5 folds (the best of them, XGBoost: \num{\W{Two}{AltAcc}}). It also retrains with about \W{Two}{CpuRatio} times less CPU time than XGBoost.
\item \textbf{Two sources of the win.} Our engine already runs a standard boosted-tree model with far less CPU than the Python libraries do:
the starting point beats XGBoost on its own. Evolution then found settings that are both more accurate and cheaper still.
\item \textbf{Tuning the libraries does not close the gap.} Searching 40 settings for XGBoost or LightGBM inside every fold gives scores near
\num{\W{Two}{XGBTunedMean}}: the extra training time costs more than the accuracy gained.
\end{itemize}

\begin{carebox}
\begin{itemize}
\item \textbf{The CPU price is part of the case.} It was fixed before any run, not taken from a real tariff. Our best program is also the most
accurate of the fully cross-validated models, so it stays ahead whenever computer time costs anything at all. Our typical run's error
(\num{\W{Two}{OursMedianRmse}}) is about the same as XGBoost's (\num{\W{Two}{AltAcc}}); its lead comes from the CPU saving.
\item \textbf{The most accurate model of all was much costlier.} On the one fold where AutoGluon was run, its 5-minute ensemble had a
somewhat lower error (\num{\W{Two}{AutoGluonFoldOneAcc}} vs our \num{\W{Two}{OursFoldOneRmse}}), but it used about
\num[round-precision=2]{\W{Two}{AutoGluonFoldOneCpu}} CPU-seconds per refit.
\item \textbf{The folds are shuffled hours.} Like every model here, ours was checked on randomly chosen hours. That measures how well
a model fits the period, not how it forecasts a future it has not seen.
\item \textbf{The other 13 cases.} We do not claim them. Under the same cross-validation either an alternative won, or our lead was within the
normal variation between folds.
\end{itemize}
\end{carebox}

\vfill
{\scriptsize\color{alt}\sloppy\textbf{Behind the numbers.} Our runs: \texttt{extend\_case.sh -{}-cuts 5}, run stamp \texttt{\BriefCvStamp}, case 02.
Alternatives: \texttt{summary\_cv.csv}, scored by \texttt{jopedime fitness evaluate}.
Built by \texttt{make -C docs/case-study-brief} from \BriefInputCount{} evidence files (digest \texttt{\BriefInputsSha}). The full technical
record of all 14 cases is \texttt{campaign-technical.pdf}.}

\end{document}
